% TOP_PROBLEM: {"id": 30006751, "title": "Asymptotic tightness of the binary Gilbert–Varshamov bound", "category_id": 15, "category": {"id": 15, "name": "computer_science", "display_name": "Computer Science", "description": "Computational complexity, algorithms, and theoretical CS.", "slug": "computer-science", "order_index": 15, "created_at": "2026-07-31T15:26:25.671Z"}, "status": "open"}
% FIRST_POSTED: null
% !TeX program = lualatex
% Compile twice: lualatex 30006751.tex
% All references and prose are embedded; no BibTeX database or external inputs.
\documentclass[11pt,a4paper]{article}
\usepackage[margin=24mm,headheight=14pt]{geometry}
\usepackage{fontspec}
\setmainfont{Latin Modern Roman}
\setsansfont{Latin Modern Sans}
\setmonofont{Latin Modern Mono}
\usepackage{amsmath,amssymb,mathtools}
\usepackage{unicode-math}
\setmathfont{Latin Modern Math}
\usepackage{microtype}
\usepackage{enumitem,xurl,needspace}
\usepackage[dvipsnames]{xcolor}
\usepackage{hyperref}
\hypersetup{unicode=true,colorlinks=true,linkcolor=MidnightBlue,urlcolor=MidnightBlue,
 pdftitle={Research attempt: Problem 30006751},pdfauthor={Research notebook},bookmarksnumbered=true}
\usepackage{bookmark}
\usepackage{tocloft}
\setlength{\cftsecnumwidth}{3em}
\cftsetpnumwidth{2.5em}
\cftsetrmarg{4em}
\usepackage{titlesec}
\titleformat{\section}{\Large\bfseries\raggedright}{\thesection}{0.7em}{}
\usepackage{fancyhdr}
\pagestyle{fancy}\fancyhf{}
\fancyhead[L]{\small\sffamily Open Problems Math | Research notebook}
\fancyhead[R]{\small Problem 30006751 | 27 September 2026}
\fancyfoot[C]{\thepage}
\setlength{\parindent}{0pt}
\setlength{\parskip}{3pt plus 1pt}
\setlength{\emergencystretch}{3em}
\setcounter{tocdepth}{1}
\setcounter{secnumdepth}{1}
\setlist[itemize]{leftmargin=3em,itemsep=5pt,topsep=4pt,parsep=0pt}
\allowdisplaybreaks
\newcommand{\N}{\mathbb{N}}
\newcommand{\Z}{\mathbb{Z}}
\newcommand{\Q}{\mathbb{Q}}
\newcommand{\R}{\mathbb{R}}
\newcommand{\C}{\mathbb{C}}
\newcommand{\F}{\mathbb{F}}
\newcommand{\A}{\mathbb{A}}
\newcommand{\PP}{\mathbb P}
\newcommand{\E}{\mathbb{E}}
\newcommand{\eps}{\varepsilon}
\newcommand{\dd}{\,\mathrm d}
\newcommand{\sref}[2]{\hyperlink{source:#1:#2}{[S#2]}}
\newcommand{\eref}[2]{\hyperlink{extra:#1:#2}{[E#2]}}
% Turn the next switch off for a shorter reading copy. The quotations preserve
% every exactTarget field, including qualifications omitted from a short summary.
\newif\ifshowcatalogue
\showcataloguetrue
\newcommand{\cataloguescope}[1]{\ifshowcatalogue
 \paragraph{Exact scope in the supplied catalogue.}
 \begin{quote}\small #1\end{quote}\fi}
\usepackage{etoolbox}
\AtBeginEnvironment{itemize}{\small}
\numberwithin{equation}{section}
% Standalone export. Original research content preserved.
\begin{document}
\setcounter{section}{0}
% ================================================================
% problemId: 30006751
% Canonical problem ID: 30006751
% exactTarget SHA-256: 9e609254f20eed4ca6127d03e9fc92dce87d19bfdae5c6fc94882834bbbcdea6
\clearpage
\section*{\texorpdfstring{Asymptotic tightness of the binary Gilbert–Varshamov bound}{Asymptotic tightness of the binary Gilbert–Varshamov bound}}\label{30006751}
\subsection{Definitions and mathematical statement}
Hamming distance counts coordinates at which binary strings differ. Define
\[
 A_2(n,d)=\max\{|C|:C\subseteq\{0,1\}^n,\ d_H(x,y)\ge d\ (x\ne y)\},
\]
\[
 R_2(\delta)=\limsup_{n\to\infty}\frac{\log_2A_2(n,\lfloor\delta n\rfloor)}n,
 \quad H_2(\delta)=-\delta\log_2\delta-(1-\delta)\log_2(1-\delta).
\]
Determine whether $R_2(\delta)=1-H_2(\delta)$ for every $0<\delta<1/2$. Codes need not be linear. The limsup, fixed relative-distance parameter, and binary alphabet are part of the assertion.
\par\smallskip\textit{Formulation and concept references:} \sref{30006751}{1}.
\cataloguescope{For integers n>=1 and d>=0, let A\_2(n,d) be the maximum cardinality of a binary code C subset \{0,1\}\textasciicircum{}n, not necessarily linear, with Hamming distance at least d between distinct codewords. For each fixed real 0<delta<1/2 define R\_2(delta)=limsup\_\{n->infinity\} n\textasciicircum{}(-1) log\_2 A\_2(n,floor(delta n)) and H\_2(delta)=-delta log\_2(delta)-(1-delta)log\_2(1-delta). Is R\_2(delta)=1-H\_2(delta) for every such delta?}
\subsection{Short English statement}
Is the binary Gilbert--Varshamov rate bound exactly optimal for every relative distance below one half, even for nonlinear codes?
% BEGIN RESEARCH ADDITION 145
\subsection{Research attempt}

\paragraph{Current frontier.}
The target permits nonlinear codes. The hierarchy completeness theorem in
the source concerns linear codes; its unrestricted version does not improve
the Delsarte bound. \sref{30006751}{1} A 2026 manuscript supplies a $\sqrt n$
multiplicative improvement to a linear-code GV existence bound, but
$\log_2(\sqrt n)/n\to0$, so that gain does not produce a strict rate
improvement. \eref{30006751}{1} Recent trace-of-AG constructions study
alphabet-reduction losses, not binary GV optimality or its failure.
\eref{30006751}{2}

\paragraph{Attack route.}
Try to disprove equality by combining a finite, possibly nonlinear binary
inner code with a large-alphabet algebraic-geometric outer family. This
would convert one finite certificate into an asymptotic counterexample.
The missing ingredient is an inner code surviving the genus penalty and
distance loss. An upper-bound route, in contrast, must remain valid for
arbitrary codes, without importing linear-code constraints.

\paragraph{Attempt: an exact finite-certificate test.}
For even $k\ge2$, let a binary inner code have length $m$, size $2^k$ and
minimum distance $d>0$, and label its words by $\F_{2^k}$. The square-field
TVZ outer guarantee is
\[
 R_{\rm out}\ge1-\Delta-\frac1{2^{k/2}-1}.
\]
The curve background is recalled in \eref{30006751}{2}. To track the loss directly,
a curve family with $N/g\to\sqrt q-1$ has evaluation codes from pole
divisors of degree $D<N$ with dimension at least $D+1-g$ and distance at
least $N-D$. Taking $D/N$ appropriately proves the displayed rate with
arbitrarily small asymptotic slack.

Coordinatewise concatenation gives length $mN$ and distance at least
$d\Delta N$, without requiring inner linearity. Therefore the usual
minimum-distance calculation certifies
\begin{equation}\label{30006751:line}
 R_2(\delta)\ge a-b\delta,\qquad
 a=\frac{k}{m}\left(1-\frac1{2^{k/2}-1}\right),\quad b=\frac{k}{d},
\end{equation}
where the line is positive. A strict gap persists after small slack in the
outer relative distance and handles all floors in the definition. The
unbounded subsequence of lengths $mN$ suffices for a limsup.

For $G(\delta)=a-b\delta-1+H_2(\delta)$, strict concavity and
$H_2'(\delta)=\log_2((1-\delta)/\delta)$ give
\begin{equation}\label{30006751:test}
 \delta_*=(1+2^b)^{-1},\qquad
 \max_{0<\delta<1/2}G(\delta)=a-1+\log_2(1+2^{-b}).
\end{equation}
Thus a sufficient finite certificate for a counterexample to the exact
asymptotic target is
\[
 \boxed{a+\log_2(1+2^{-b})>1.}
\]
If it holds, the line at $\delta_*$ is positive automatically because
$1-H_2(\delta_*)>0$, so the outer positive-rate range is respected. A
negative value means this certified line fails, not that the concatenated
codes could not have better distances than the product lower bound.

\paragraph{One entire inner family fails the certificate.}
Take all affine functions on $\F_2^{k-1}$ as words. Then $m=2^{k-1}$,
size $2^k$, and $d=m/2$. Write $\rho=k/m$ and $Q=2^{k/2}$. The maximum
gap in \eqref{30006751:test} equals
\[
 -\frac{\rho}{Q-1}+\log_2\cosh(\rho\log2)
 \le-\frac{\rho}{Q-1}+\frac{\rho^2\log2}{2}<0.
\]
The inequality $\log\cosh x\le x^2/2$ follows by integrating
$\tanh x\le x$. To verify strict negativity, put $k=2j$ and use
$2j\le2^j$ for $j\ge1$, proved by induction starting at $j=1,2$.
Then $k(Q-1)<Q^2=2^k$, so $\rho(Q-1)\log2<2$. This proves the failure
uniformly for every even $k$. It excludes only the TVZ/product-distance
calculation, not an analysis retaining the full distance distributions.

\paragraph{A nonlinear inner-code experiment with an exact construction.}
Use coordinates $x\in\F_2^4$. Order the pairs as
$(1,2),(1,3),(1,4),(2,3),(2,4),(3,4)$; the six binary digits of $z$ are the
coefficients of $Q_z(x)=\sum_{i<j}z_{ij}x_ix_j$. Set
\[
 Z=\{0,12,22,27,39,41,53,58\},
\]
\[
 C=\{(Q_z(x)+\ell\cdot x+c)_{x\in\F_2^4}:
       z\in Z,\ \ell\in\F_2^4,\ c\in\F_2\}.
\]
All $\binom82$ differences have nondegenerate polar forms of rank four,
verified exactly by the companion script. Squaring a quadratic Walsh sum
and translating by $h$ leaves only $h=0$, because every nonzero polar linear
form sums to zero. The Walsh sums are therefore $\pm4$. Different quadratic
cosets have distances $6$ or $10$; within a coset distinct affine words have
distance $8$ or $16$. The cosets are disjoint, producing a $(16,256,6)$ code.
Its quadratic parts are not closed under addition, so it is nonlinear.
These familiar finite parameters are reconstructed here without a novelty
claim.

The script also enumerates the binary Hamming code, its parity extension,
and the parity extension of the length-$23$ cyclic code generated by
$X^{11}+X^9+X^7+X^6+X^5+X+1$. Exact minimum distances and the corresponding
values of \eqref{30006751:test} are
\[
\begin{array}{c|c|c|r}
 m&\log_2|C|&d&\max G\ \text{(approximately)}\\\hline
 7&4&3&-0.13687024\\
 8&4&4&-0.08170417\\
 16&8&6&-0.05115596\\
 24&12&8&-0.07118471
\end{array}
\]
The signs do not depend on floating-point logarithms. For length $16$,
$2^{-4/3}<2/5$ and $(7/5)^{15}<2^8$, so the logarithm is below
$8/15=1-a$. The same certificate works at length $7$, where
$1-a=13/21>8/15$. At length $8$, $(3/2)^3<2^2$. At length $24$,
$2^{-3/2}<9/25$, $(34/25)^2<2$, and $1-a=32/63>1/2$.

\paragraph{Stress test.}
Cartesian powers of a fixed code do not retain relative minimum distance:
two product words can differ in one block, giving distance $d$ and relative
distance $d/(mN)\to0$. The outer code is essential. Polynomial and other
$\exp(o(n))$ size improvements disappear in normalized logarithmic rate.
Neither the finite failures nor the infinite affine-family calculation
excludes other inner codes, multilevel methods or distance-spectrum gains.
Conversely, linear-hierarchy completeness cannot establish a nonlinear
upper bound by silently assuming linear-code identities.

\paragraph{Outcome.}
\textbf{Outcome: Exploratory approach with unresolved gap.}
The attempt derives a sufficient finite-code test for an asymptotic
counterexample and proves that the Hadamard/TVZ parameter family never
passes it. Four finite codes, including an explicitly checked nonlinear
one, also fail the test. No passing certificate and no matching upper bound
for arbitrary binary codes have been obtained.

% END RESEARCH ADDITION 145
\subsection{Sources}
\begin{itemize}\raggedright
\item[\textnormal{[S1]}]\hypertarget{source:30006751:1}{} Leonardo Nagami Coregliano, Fernando Granha Jeronimo and Chris Jones, A Complete Linear Programming Hierarchy for Linear Codes; ITCS2022 article51.
\url{https://drops.dagstuhl.de/storage/00lipics/lipics-vol215-itcs2022/LIPIcs.ITCS.2022.51/LIPIcs.ITCS.2022.51.pdf}.
{\small\textit{Catalogue formulation source.}}
% BEGIN RESEARCH REFERENCES 145
\item[\textnormal{[E1]}]\hypertarget{extra:30006751:1}{} C. Yuan and R. Zhu,
\emph{Improvement of the Gilbert--Varshamov Bound for Linear Codes and Quantum
Codes}, arXiv:2601.18590v1 (26 January 2026), Theorems 1.3 and 4.1.
\url{https://arxiv.org/html/2601.18590v1}.
\item[\textnormal{[E2]}]\hypertarget{extra:30006751:2}{} G. Cohen, D. Doron,
N. Goldgraber and T. Manket, \emph{Tracing AG Codes: Toward Meeting the
Gilbert--Varshamov Bound}, ICALP 2026, LIPIcs 374, article 68, Sections 1,
2.1 and 8 (optimal curves and alphabet-reduction limitations).
\url{https://drops.dagstuhl.de/storage/00lipics/lipics-vol374-icalp2026/html/LIPIcs.ICALP.2026.68/LIPIcs.ICALP.2026.68.html}.

% END RESEARCH REFERENCES 145
\end{itemize}

\end{document}
